\chapter{Einleitung}\label{ch:einleitung}

\section{Was ist FourierGrating?}
\FG{} berechnet, wie eine elektromagnetische Welle an einer \emph{periodisch strukturierten} Oberfläche oder
einem Schichtstapel gebeugt wird. Typische Ergebnisse sind
\begin{itemize}
  \item die \textbf{Beugungseffizienzen} aller reflektierten und transmittierten Ordnungen, getrennt nach s- und
        p-Anteil, mit komplexen Amplituden,
  \item die \textbf{Absorption} und die Energiebilanz,
  \item \textbf{Nahfelder} ($\mathbf E$, $\mathbf H$, $|\mathbf E|^2$, Poynting-Vektor) in Schnitten durch die Struktur,
  \item bei isolierten Strukturen \textbf{Streubreiten und Streuquerschnitte}, bei Wellenleitern die Leistung
        geführter \textbf{Moden},
  \item \textbf{Gradienten} beliebiger Zielfunktionen der Effizienzen nach Geometrie-, Material- und
        Einfallsparametern, für die Optimierung.
\end{itemize}

Das Programm löst die Maxwellgleichungen mit der \textbf{Fourier-Modal-Methode (FMM)}. Die Felder werden in den
periodischen Richtungen in Fourierreihen entwickelt; in der Schichtrichtung $y$ entsteht ein System gewöhnlicher
Differentialgleichungen $\mathrm dF/\mathrm dy = M(y)\,F$. \FG{} bietet zwei Wege, es zu lösen:

\begin{description}
  \item[RCWA] (\emph{Rigorous Coupled-Wave Analysis}): Die Struktur wird in Scheiben zerlegt, in denen $M$ nicht von
        $y$ abhängt. In jeder Scheibe wird ein Eigenwertproblem gelöst. Das ist ideal für Strukturen mit senkrechten
        Wänden (Lamellen, Säulen, Löcher).
  \item[Differentialmethode] mit \textbf{implizitem Runge--Kutta-Verfahren} (Gauß--Legendre, Radau~IIA): $M(y)$ wird
        direkt integriert. Das ist ideal für Profile mit schrägen oder gekrümmten Flanken (Sinus- und
        Echelettegitter), weil keine Treppenstufen entstehen.
\end{description}
Beide Wege verwenden die \textbf{schnelle Fourierfaktorisierung} (\emph{Fast Fourier Factorization}, FFF) nach
Nevière und Popov mit Normalenvektorfeldern. Sie sorgt dafür, dass auch Metallgitter in TM-Polarisation schnell
konvergieren.

\section{Funktionsumfang im Überblick}
\begin{center}
\small
\begin{tabularx}{\linewidth}{>{\raggedright\arraybackslash}p{4.1cm}X}
\toprule
Bereich & Möglichkeiten \\
\midrule
Geometrie & 1D-Gitter (periodisch in $x$), 2D-Gitter (periodisch in $x$ und $z$), isolierte Strukturen mit PML,
rotationssymmetrische Strukturen (in Arbeit) \\
Schichttypen & homogen, Lamellen, Profile $y=g(x)$ (Sinus, Dreieck, Spline), Polygone, gekrümmte Konturen (Kreis,
Ellipse, NURBS), Profile $y=g(x,z)$, STEP-Körper aus CAD, Ringe \\
Materialien & isotrop (reell oder komplex), anisotrop mit vollem $\varepsilon$-Tensor oder Hauptwerten und
Euler-Winkeln; Metalle, Dielektrika, Halbleiter \\
Einfall & ebene Welle beliebiger Richtung und Polarisation (auch konisch), Gauß-Strahl (Zylinderpfad), geführte Mode
eines strukturierten Halbraums \\
Löser & RCWA und Runge--Kutta; Faktorisierungen Laurent, Li, FFF in mehreren Varianten; adaptive räumliche Auflösung
(ASR) \\
Auswertung & Effizienzen, s/p-Anteile, Amplituden, S-Matrix in der s/p-Basis, Nahfelder, Streuquerschnitte,
Modenleistungen \\
Scans & Wellenlänge, Winkel und Polarisation als Liste oder Bereich (kartesisches Produkt); Sweep über beliebige Zahlen
der Konfiguration \\
Optimierung & exakte Gradienten über die diskrete Adjungierte; Adapter für PyTorch und JAX \\
Leistung & OpenMP-Parallelisierung auf drei Ebenen, CUDA-Backend für NVIDIA-GPUs, wählbare Genauigkeit
(fp64/mixed/fp32) \\
Schnittstellen & Kommandozeile \cmd{fgsolve}, Python-Paket \texttt{fouriergrating}, C++-Bibliothek, grafische
Oberfläche im Browser (\GUI) \\
\bottomrule
\end{tabularx}
\end{center}

\section{Was FourierGrating (noch) nicht kann}
\begin{itemize}
  \item \textbf{Dispersive Materialien}: Ein Material hat für alle Wellenlängen eines Scans dieselbe Brechzahl.
        Für Spektren über einen größeren Bereich ist das bei Metallen und Halbleitern falsch. Abhilfe bis zur
        Umsetzung (Meilenstein M22, geplant): je Wellenlänge eine Rechnung mit passender Brechzahl, am bequemsten in
        Python (Abschnitt~\ref{sec:py-dispersion}).
  \item Nichtlineare Optik, magnetische Materialien ($\mu_r\neq1$), zeitabhängige Felder.
  \item Kurvenförmige Schichtgrenzen in $y$-Richtung außerhalb der unterstützten Profil- und CAD-Typen.
  \item Rotationssymmetrische Strukturen (Meilenstein M20) sind teilweise umgesetzt; Nahfeld, Profile in Ringform
        und CAD-Rotationskörper folgen. $P$-zählige Drehsymmetrie (M21) ist spezifiziert, aber noch nicht umgesetzt.
\end{itemize}

\section{Aufbau dieses Handbuchs}
\begin{center}
\small
\begin{tabularx}{\linewidth}{lY}
\toprule
Kapitel & Inhalt \\
\midrule
\ref{ch:grundlagen} & Physikalische Grundlagen, Koordinaten, Definitionen -- \emph{einmal lesen} \\
\ref{ch:installation}, \ref{ch:schnellstart} & Installation und erste Rechnung \\
\ref{ch:konfiguration}--\ref{ch:ausgabe} & Konfigurationsdatei, Schichttypen, Löserwahl, Scans, Nahfeld, Ausgaben
-- \emph{das Nachschlagewerk} \\
\ref{ch:erweitert}, \ref{ch:cad} & Isolierte Strukturen, ASR, Zylinderkoordinaten, CAD-Import \\
\ref{ch:python}, \ref{ch:optimierung}, \ref{ch:cpp} & Programmierschnittstellen und Optimierung \\
\ref{ch:leistung} & Rechenzeit, Parallelisierung, GPU, Genauigkeit \\
\ref{ch:gui} & Grafische Oberfläche im Browser \\
\ref{ch:rezepte} & Durchgerechnete Beispiele \\
\ref{ch:probleme} & Fehlermeldungen und Fehlersuche \\
Anhang & Schlüsselreferenz, Umgebungsvariablen, Referenzfälle, Glossar, Literatur \\
\bottomrule
\end{tabularx}
\end{center}

\section{Darstellung}
\begin{itemize}
  \item Schlüssel der Konfigurationsdatei erscheinen als \key{lambda\_um}, Werte als \val{"rcwa"}, Dateien als
        \file{examples/v3\_lamellar\_dielectric.json}, Befehle als \cmd{fgsolve}.
  \item Pfade sind relativ zum Wurzelverzeichnis des Repositorys \file{FourierGrating}, sofern nicht anders gesagt.
  \item Einheiten: Längen immer in Mikrometern (\textmu m), Winkel in Grad. Das gilt für die Konfiguration, die
        Python-Schnittstelle und die GUI gleichermaßen.
  \item Kommentare mit \texttt{//} in JSON-Beispielen dienen nur der Erklärung. Echte Konfigurationsdateien sind
        reines JSON ohne Kommentare.
\end{itemize}

\begin{hinweis}[Verbindliche Quelle]
Alle Formeln, Algorithmen und Abnahmekriterien stehen im \textbf{Implementierungsplan}
(\file{docs/implementation\_plan/main.pdf}). Dieses Handbuch erklärt die \emph{Benutzung}. Wo es Formeln nennt,
verweist es auf den Plan; bei einem Widerspruch gilt der Plan bzw.\ das tatsächliche Verhalten des Programms.
\end{hinweis}
